\documentclass[10pt,a4paper]{amsart}
\usepackage{amsmath,amssymb,amsthm,graphicx}
\usepackage{verbatim}
\newenvironment{reference}{\comment}{\endcomment}
%\newenvironment{reference}{}{}
\newenvironment{slogan}{\comment}{\endcomment}
\newenvironment{history}{\comment}{\endcomment}

% For commutative diagrams we use Xy-pic
\usepackage[all]{xy}

% We use 2cell for 2-commutative diagrams.
\xyoption{2cell}
% Theorem environments.
%
\theoremstyle{plain}
\newtheorem{theorem}[subsection]{Theorem}
\newtheorem{proposition}[subsection]{Proposition}
\newtheorem{lemma}[subsection]{Lemma}

\theoremstyle{definition}
\newtheorem{definition}[subsection]{Definition}
\newtheorem{example}[subsection]{Example}
\newtheorem{exercise}[subsection]{Exercise}
\newtheorem{situation}[subsection]{Situation}

\theoremstyle{remark}
\newtheorem{remark}[subsection]{Remark}
\newtheorem{remarks}[subsection]{Remarks}

\numberwithin{equation}{subsection}

% Macros
%
\def\lim{\mathop{\mathrm{lim}}\nolimits}
\def\colim{\mathop{\mathrm{colim}}\nolimits}
\def\Spec{\mathop{\mathrm{Spec}}}
\def\Hom{\mathop{\mathrm{Hom}}\nolimits}
\def\Ext{\mathop{\mathrm{Ext}}\nolimits}
\def\SheafHom{\mathop{\mathcal{H}\!\mathit{om}}\nolimits}
\def\SheafExt{\mathop{\mathcal{E}\!\mathit{xt}}\nolimits}
\def\Sch{\mathit{Sch}}
\def\Mor{\mathop{\mathrm{Mor}}\nolimits}
\def\Ob{\mathop{\mathrm{Ob}}\nolimits}
\def\Sh{\mathop{\mathit{Sh}}\nolimits}
\def\NL{\mathop{N\!L}\nolimits}
\def\CH{\mathop{\mathrm{CH}}\nolimits}
\def\proetale{{pro\text{-}\acute{e}tale}}
\def\etale{{\acute{e}tale}}
\def\QCoh{\mathit{QCoh}}
\def\Ker{\mathop{\mathrm{Ker}}}
\def\Im{\mathop{\mathrm{Im}}}
\def\Coker{\mathop{\mathrm{Coker}}}
\def\Coim{\mathop{\mathrm{Coim}}}

% Boxtimes
%
\DeclareMathSymbol{\boxtimes}{\mathbin}{AMSa}{"02}

%
% Macros for moduli stacks/spaces
%
\def\QCohstack{\mathcal{QC}\!\mathit{oh}}
\def\Cohstack{\mathcal{C}\!\mathit{oh}}
\def\Spacesstack{\mathcal{S}\!\mathit{paces}}
\def\Quotfunctor{\mathrm{Quot}}
\def\Hilbfunctor{\mathrm{Hilb}}
\def\Curvesstack{\mathcal{C}\!\mathit{urves}}
\def\Polarizedstack{\mathcal{P}\!\mathit{olarized}}
\def\Complexesstack{\mathcal{C}\!\mathit{omplexes}}
% \Pic is the operator that assigns to X its picard group, usage \Pic(X)
% \Picardstack_{X/B} denotes the Picard stack of X over B
% \Picardfunctor_{X/B} denotes the Picard functor of X over B
\def\Pic{\mathop{\mathrm{Pic}}\nolimits}
\def\Picardstack{\mathcal{P}\!\mathit{ic}}
\def\Picardfunctor{\mathrm{Pic}}
\def\Deformationcategory{\mathcal{D}\!\mathit{ef}}
\usepackage{cleveref}
\crefname{originalreference}{original reference}{original references}
\Crefname{originalreference}{Original reference}{Original references}
\makeatletter
\newlabel{lemma-count-embeddings-explicitly}{{09H8}{source}{}{sourcepointer-1}{}}
\newlabel{lemma-count-embeddings-explicitly@cref}{{[originalreference][0][]09H8}{[1][0][]source}{}{}{}}
\makeatother
\crefname{originalreference}{original reference}{original references}
\Crefname{originalreference}{Original reference}{Original references}
\makeatletter
\newlabel{lemma-degree-minimal-polynomial}{{09GN}{source}{}{sourcepointer-2}{}}
\newlabel{lemma-degree-minimal-polynomial@cref}{{[originalreference][0][]09GN}{[1][0][]source}{}{}{}}
\makeatother
\crefname{originalreference}{original reference}{original references}
\Crefname{originalreference}{Original reference}{Original references}
\makeatletter
\newlabel{lemma-multiplicativity-degrees}{{09G9}{source}{}{sourcepointer-3}{}}
\newlabel{lemma-multiplicativity-degrees@cref}{{[originalreference][0][]09G9}{[1][0][]source}{}{}{}}
\makeatother
\title{Embedding count equals degree exactly when each step of the generating tower is separable}
\author{The Stacks Project Authors}
\date{stacks-project:fields.tex:local-snapshot-20261012; GNU FDL 1.2 or later}
\begin{document}
\maketitle
\noindent\textit{Editorial note.} All delivered mathematical text is the author's own, copied from the pinned Stacks Project chapter file fields.tex; no statement or proof step was rewritten, reordered or removed.\par\medskip
\noindent\textit{Editorial note.} The counted proof is the author's complete argument: The reordering trick in the last paragraph uses that the strict-inequality and equality parts of this very lemma have already been proved within the counted span; gamma\_2, ..., gamma\_m are filled by alpha\_1, ..., alpha\_n (source lines 1589-1592).\par\medskip
\noindent\textit{Editorial note.} The label lemma-separably-generated-separable names the counted statement, and the proof environment immediately following it in the source is delivered verbatim as its proof.\par\medskip
\noindent\textit{Editorial note.} The counted statement is carried by the author's own lemma environment in the chapter source.\par\medskip
\noindent\textit{Allowed standard prerequisites (editorial).}\begin{itemize}
\item Context span 1485-1496: the Situation (generating tower, minimal polynomials P\_i, fixed algebraic closure)
\item Same chapter above the span: lemma-count-embeddings-explicitly (|Mor\_F(K, F-bar)| = product deg\_s(P\_i), source lines 1542-1554), lemma-degree-minimal-polynomial (deg(P\_i) = [K\_i : K\_\{i-1\}]), lemma-multiplicativity-degrees (degrees multiply in towers)
\end{itemize}\medskip
\ifdefined\hypertarget\hypertarget{sourcepointer-1}{}\fi
\noindent\textit{Source reference.} Original reference 09H8: sources/stacks-fields/pinned-original-source.tex, source lines 1543--1543. Prerequisite pointer invoked at 1581; it names a standard fact of the same chapter that the counted argument uses, and every step of the counted proof built on it is retained verbatim in the delivered body.\par\medskip
\ifdefined\hypertarget\hypertarget{sourcepointer-2}{}\fi
\noindent\textit{Source reference.} Original reference 09GN: sources/stacks-fields/pinned-original-source.tex, source lines 1087--1087. Prerequisite pointer invoked at 1575; it names a standard fact of the same chapter that the counted argument uses, and every step of the counted proof built on it is retained verbatim in the delivered body.\par\medskip
\ifdefined\hypertarget\hypertarget{sourcepointer-3}{}\fi
\noindent\textit{Source reference.} Original reference 09G9: sources/stacks-fields/pinned-original-source.tex, source lines 759--759. Prerequisite pointer invoked at 1576; it names a standard fact of the same chapter that the counted argument uses, and every step of the counted proof built on it is retained verbatim in the delivered body.\par\medskip
\begin{situation}
\label{situation-finitely-generated}
Here $F$ be a field and $K/F$ is a finite extension generated by elements
$\alpha_1, \ldots, \alpha_n \in K$. We set $K_0 = F$ and
$$
K_i = F(\alpha_1, \ldots, \alpha_i)
$$
to obtain a tower of finite extensions
$K = K_n / K_{n - 1} / \ldots / K_0 = F$.
Denote $P_i$ the minimal polynomial of $\alpha_i$ over $K_{i - 1}$.
Finally, we fix an algebraic closure $\overline{F}$ of $F$.
\end{situation}

\begin{lemma}
\label{lemma-separably-generated-separable}
Assumptions and notation as in Situation \ref{situation-finitely-generated}.
If each $P_i$ is separable, i.e., each $\alpha_i$ is separable over
$K_{i - 1}$, then
$$
|\Mor_F(K, \overline{F})| = [K : F]
$$
and the field extension $K/F$ is separable. If one of the $\alpha_i$ is
not separable over $K_{i - 1}$, then
$|\Mor_F(K, \overline{F})| < [K : F]$.
\end{lemma}

\begin{proof}
If $\alpha_i$ is separable over $K_{i - 1}$ then
$\deg_s(P_i) = \deg(P_i) = [K_i : K_{i - 1}]$
(last equality by Lemma \ref{lemma-degree-minimal-polynomial}).
By multiplicativity (Lemma \ref{lemma-multiplicativity-degrees}) we have
$$
[K : F] = \prod [K_i : K_{i - 1}] = \prod \deg(P_i) =
\prod \deg_s(P_i) = |\Mor_F(K, \overline{F})|
$$
where the last equality is Lemma \ref{lemma-count-embeddings-explicitly}.
By the exact same argument we get the strict inequality
$|\Mor_F(K, \overline{F})| < [K : F]$ if one of the $\alpha_i$ is
not separable over $K_{i - 1}$.

\medskip\noindent
Finally, assume again that each $\alpha_i$ is separable over $K_{i - 1}$.
We will show $K/F$ is separable.
Let $\gamma = \gamma_1 \in K$ be arbitrary. Then we can find additional
elements $\gamma_2, \ldots, \gamma_m$ such that
$K = F(\gamma_1, \ldots, \gamma_m)$ (for example we could take
$\gamma_2 = \alpha_1, \ldots, \gamma_{n + 1} = \alpha_n$).
Then we see by the last part of the lemma (already proven above)
that if $\gamma$ is not separable over $F$ we would have the
strict inequality $|\Mor_F(K, \overline{F})| < [K : F]$
contradicting the very first part of the lemma (already prove above
as well).
\end{proof}

\end{document}
